\title{xLSTM: Extended Long Short-Term Memory}

\begin{abstract}
The foundational innovations of the Long Short-Term Memory (LSTM)—specifically, the constant error carousel and gating mechanisms—were introduced in the 1990s. Since their inception, LSTMs have proven highly durable and have played a key role in the advancement of deep learning, serving as the basis for early Large Language Models (LLMs). Nevertheless, the rise of Transformer architectures, characterized by their parallelizable self-attention mechanisms, signaled a transformative shift and enabled surpassing LSTMs in large-scale tasks. This work revisits an essential inquiry: What potential does language modeling have when LSTMs are scaled to billions of parameters, employ the recent advances developed for modern LLMs, and address persistent weaknesses of LSTMs? Our first contribution is the introduction of exponential gating, augmented by targeted normalization and stabilization strategies. Secondly, we redesign the LSTM memory, resulting in: (i) sLSTM, which incorporates scalar memory, scalar updates, and a revised approach to memory mixing, and (ii) mLSTM, which is fully parallelizable, utilizes a matrix memory, and employs a covariance-based update rule. By embedding these LSTM variants within residual block frameworks, we create xLSTM blocks, which are then sequentially composed to form xLSTM architectures. The integration of exponential gating and redesigned memory significantly enhances the expressiveness and performance of xLSTM models, allowing them to achieve competitive or superior results compared to cutting-edge Transformers and State Space Models, especially in terms of scalability and overall capability.\\
Code available at: \url{https://github.com/NX-AI/xlstm}
\end{abstract}

\section{Introduction}
\label{sec:intro}

The Long Short-Term Memory (LSTM) ideas~\citep{Hochreiter:91,Hochreiter:97nips,Hochreiter:97}, i.e., the constant error carousel and gating, were introduced to overcome the vanishing gradient problem of recurrent neural networks~\citep{Hochreiter:91,Hochreiter:00book}:

\begin{align}
\label{eq:lstm_idea}
  c_t \ = \ f_t \ c_{t-1} \ + \ 
          i_t \  z_t \ , \quad  
          h_t \ = \ o_t \ \psi({c_t}) \ . 
\end{align}

The constant error carousel operates by incrementally adjusting the cell state $c_{t-1}$ (depicted in green) through the addition of cell inputs $z_t$, with this process being regulated by sigmoid gates (shown in blue). The input gate $i_t$ alongside the forget gate $f_t$ govern the modifications to the cell state, whereas the output gate $o_t$ determines the final output of the memory cell, which corresponds to the hidden state $h_t$. The cell state undergoes normalization or squashing via $\psi$ prior to the application of the output gate, which ultimately yields the hidden state.

LSTMs have demonstrated broad applicability across numerous domains \citep{Hochreiter:01,Hochreiter:07,Schmidhuber:15}, dominating text generation approaches prior to the emergence of Transformers in 2017~\citep{Vaswani:17}. Their capabilities have been proven in a variety of sequential tasks, such as text generation~\citep{Graves:13,Karpathy:15blog}, handwriting synthesis~\citep{Graves:13}, sequence-to-sequence machine translation~\citep{Sutskever:14nips}, program evaluation~\citep{Zaremba:14arxiv}, image captioning~\citep{Karpathy:15,Hossain:19}, code generation~\citep{Karpathy:15blog}, hydrological forecasting including rainfall-runoff modeling~\citep{Kratzert:18,Kratzert:19}, and in the development of flood warning systems~\citep{Nearing:24}. Within the realm of reinforcement learning, LSTMs are among the most effective sequence models, forming the backbone of agents such as AlphaStar for StarCraft II~\citep{Vinyals:17short}, OpenAI Five for Dota 2~\citep{OpenAI:19}, as well as controllers for magnetic systems in nuclear fusion~\citep{Degrave:22short}. LSTMs are particularly well-suited for abstracting high-level features, efficiently encoding semantic concepts within their memory cells~\citep{Karpathy:15blog}. This proficiency has been empirically observed through the discovery of number and syntax neurons~\citep{Lakretz:19}, linguistic-specific neurons~\citep{Bau:19}, and sentiment neurons~\citep{Radford:17arxiv}. Despite advancements in alternative architectures, LSTMs remain prevalent in critical applications~\citep{Degrave:22short,Nearing:24} and continue to demonstrate lasting relevance.

LSTMs, though highly effective in a variety of tasks, are constrained by three principal shortcomings: (i) Difficulty in updating memory allocations. This challenge is illustrated by the 
{\it Nearest Neighbor Search} problem (for further discussion, see Appendix~\ref{sec:appNearestNeighborSearch}): given a reference vector, the model must iterate through a sequence to identify and recall the most similar vector for subsequent value retrieval. As displayed in the left panel of Figure~\ref{fig:lstmProblems}, LSTM networks exhibit notable mean squared error on this task, reflecting their struggle to overwrite previously stored information even when a superior match is encountered later in the sequence. In contrast, the xLSTM, introduced here, addresses this difficulty through the use of exponential gating mechanisms. (ii) Constrained memory capacity, whereby all information is required to be condensed into scalar-valued cell states. This bottleneck is made evident by the {\it Rare Token Prediction} experiment. In the right panel of Figure~\ref{fig:lstmProblems}, token prediction perplexity on Wikitext-103~\citep{Merity:17} is broken down by the frequency of token occurrence. The standard LSTM demonstrates poorer performance for infrequent tokens due to its restrictive memory representation; the xLSTM alleviates this issue by incorporating a matrix-based memory. (iii) Inefficiency for parallel computation caused by recurrent connections—specifically, the hidden-to-hidden transitions that necessitate strictly sequential data processing.

The above limitations of LSTMs have motivated the widespread adoption of Transformers~\citep{Vaswani:17} for language modeling tasks. This naturally raises the question: how well can language modeling tasks be addressed by LSTMs if these shortcomings are overcome and the models are scaled to match the parameter counts of modern Large Language Models?

\section{Extended Long Short-Term Memory}
\label{sec:xLSTM}

In addressing the shortcomings of traditional LSTM, the Extended Long Short-Term Memory (xLSTM) proposes two principal advancements to the LSTM framework described in Equation~\eqref{eq:lstm_idea}. These advancements—specifically, exponential gating and innovative memory mechanisms—expand the LSTM landscape with two distinctive variants: (i) sLSTM (discussed in Section~\ref{sec:sLSTM}), which operates with a scalar memory, a scalar-based update, and integrates memory mixing; and (ii) mLSTM (see Section~\ref{sec:mLSTM}), which leverages matrix-form memory alongside an outer product (covariance-based) update rule that is designed to be fully parallelizable. Both sLSTM and mLSTM utilize exponential gating to enhance the original LSTM capabilities. To enable efficient parallelism, mLSTM discards the conventional hidden-to-hidden recurrent (memory mixing) connections. These models may also be generalized to configurations involving multiple memory cells, with sLSTM allowing memory mixing between cells. Furthermore, the sLSTM architecture supports multiple attention heads without inter-head memory mixing, while maintaining cell-wise memory integration within each head, thus presenting a new paradigm for memory mixing enabled by both multiple heads and exponential gating. In the case of mLSTM, the utilization of multiple heads or cells yields equivalent functionality.

By incorporating these advanced LSTM forms within residual block structures, the architecture gives rise to xLSTM blocks (refer to Section~\ref{sec:xLSTMblock}). Layering these blocks residually forms comprehensive xLSTM-based neural architectures (see Section~\ref{sec:xLSTMarchitecure}). An overview and visualization of the xLSTM architecture and its building blocks can be found in Figure~\ref{fig:xlstm_sketch}.

\subsection{Review of the Long Short-Term Memory}
\label{sec:orgLSTM}

The initial formulation of LSTM~\citep{Hochreiter:91,Hochreiter:97nips,Hochreiter:97} proposed the scalar memory cell as the primary component for both computation and information retention, effectively addressing the vanishing gradient problem~\citep{Hochreiter:91,Hochreiter:00book} via the constant error carousel mechanism, realized through updates to the cell state. This cell incorporates three types of gates: input, output, and the forget gate—the latter introduced later by \citet{Gers:00}. At each discrete time point $t$, the LSTM memory cell is governed by the following update equations:

\begin{align}
c_t \ &= \  f_t \ c_{t-1} \ + \ i_t \ z_t &  & &\text{cell state} \\
h_t  \ &= \ o_t \ \tilde{h}_t \ , & \tilde{h}_t \ &= \ \psi \left( c_t\right)
  &\text{hidden state} \\
z_t \ &= \ \varphi \left( \tilde{z}_t \right) \ , 
  &\tilde{z}_t \ &=  \ \mathbf{w}^\top_{z} \ \mathbf{x}_t \ + \
  r_{z}  h_{t-1} \ + \  b_{z} \ \
  &\text{cell input} \\
i_t \ &= \ \sigma \left( \tilde{i}_t  \right) \ , 
  &\tilde{i}_t \ &= \ \mathbf{w}^\top_{i} \ \mathbf{x}_t \ + \
  r_{i}  \ h_{t-1} \ + \  b_{i} \ \
  &\text{input gate} \\
f_t \ &= \ \sigma \left(  \tilde{f}_t \right) \ , 
  &\tilde{f}_t \ &= \ \mathbf{w}^\top_{f} \ \mathbf{x}_t  \ + \
  r_{f}  \ h_{t-1} \ + \  b_{f} \ \
  &\text{forget gate} \\
o_t \ &= \ \sigma \left( \tilde{o}_t \right) \ , 
  &\tilde{o}_t  \ &= \ \mathbf{w}^\top_{o} \ \mathbf{x}_t \ + \
  r_{o}  \ h_{t-1} \ + \  b_{o} \ \
  &\text{output gate} 
\end{align}

The input weight vectors $\mathbf{w}_{z}$, $\mathbf{w}_{i}$, $\mathbf{w}_{f}$, and $\mathbf{w}_{o}$ map the input $\mathbf{x}_t$ to the respective units: cell input, input gate, forget gate, and output gate. The parameters $r_{z}$, $r_{i}$, $r_{f}$, and $r_{o}$ are the recurrent weights that connect the previous hidden state $h_{t-1}$ with each of these same units. The biases associated with each component are denoted as $b_{z}$, $b_{i}$, $b_{f}$, and $b_{o}$. The activation functions for the cell input and the hidden state are denoted by $\varphi$ and $\psi$ (with $\tanh$ commonly used in practice). The function $\psi$ specifically serves to constrain the range of the cell state, which might otherwise grow without bounds. The activation for each gate uses the sigmoid nonlinearity, i.e., $\sigma(x) = 1/(1+\exp(-x))$. More recent versions aggregate multiple scalar memory cells into a vector, $\mathbf{c}_t \in \mathbb{R}^{d}$, enabling the use of recurrent matrices $\mathbf{R} \in \mathbb{R}^{d \times d}$ that permit intermixing of outputs from different memory units~\citep{Greff:15}; see further discussion in Appendix~\ref{sec:apporgLSTM}. Ablation experiments have demonstrated that every component within the memory cell architecture plays an essential role~\citep{Greff:15}.

\subsection{sLSTM}
\label{sec:sLSTM}

To enable LSTMs to update their memory management, we incorporate exponential gates (shown in red), along with mechanisms for normalization and stabilization. Specifically, we allow input and forget gates to utilize exponential activation functions. For normalization purposes, we propose a normalizer state, which accumulates the sum of the input gate multiplied by all subsequent forget gates.

The scalar sLSTM forward pass is:

\begin{align}
\label{eq:slstmforward}
c_t \ &= \  f_t \ c_{t-1} \ + \ i_t \ z_t & & 
 &\text{cell state} \\
n_t \ &= \  f_t \ n_{t-1} \ + \ i_t  & & 
 &\text{normalizer state} \\
h_t \ &= \ o_t \ \tilde{h}_t \ ,
  &\tilde{h}_t \ &= \ c_t / n_t
&\text{hidden state} \\
z_t \ &= \ \varphi \left( \tilde{z}_t \right) \ , 
  &\tilde{z}_t \ &=  \ \mathbf{w}^\top_{z} \ \mathbf{x}_t \ + \
  r_{z}  h_{t-1} \ + \  b_{z}
  &\text{cell input} \\
i_t \ &= \ \!\exp \left( \tilde{i}_t  \right) \ , 
  &\tilde{i}_t \ &= \ \mathbf{w}^\top_{i} \ \mathbf{x}_t \ + \
  r_{i}  \ h_{t-1} \ + \  b_{i}
  &\text{input gate} \\
f_t \ &= \ \sigma \left(  \tilde{f}_t \right) \ \text{OR} \
\!\exp \left(  \tilde{f}_t \right) \ ,
  &\tilde{f}_t \ &= \ \mathbf{w}^\top_{f} \ \mathbf{x}_t  \ + \
  r_{f}  \ h_{t-1} \ + \  b_{f}
  &\text{forget gate} \\
o_t \ &= \ \sigma \left( \tilde{o}_t \right) \ , 
  &\tilde{o}_t  \ &= \ \mathbf{w}^\top_{o} \ \mathbf{x}_t \ + \
  r_{o}  \ h_{t-1} \ + \  b_{o}
  &\text{output gate} 
\end{align}

We transfer the original LSTM gating techniques, i.e., input- and/or hidden-dependent gating plus bias term, to the new architectures. Exponential activation functions can lead to large values that cause overflows. Therefore, we stabilize gates with an additional state \mbox{$m_t$~\citep{Milakov:18arxiv}}:

\begin{align}
\label{eq:slstmstabil}
m_t \ &= \ \max \left( \log ( f_t ) + m_{t-1} , \log ( i_t ) \right) &\text{stabilizer state} \\
i'_t \ &= \ \!\exp \left( \log \left ( i_t \right) - m_t \right) \ = \ \!\exp \left( \tilde{i}_t - m_t \right) \ \, 
  &\text{stabil. input gate} \\
f'_t \ &= \ \!\exp \left( \log \left( f_t \right) + m_{t-1} - m_t \right) \ \, 
  &\text{stabil. forget gate}
\end{align}

As demonstrated in Appendix~\ref{sec:appsLSTM}, substituting $f_t$ with $f'_t$ and $i_t$ with $i'_t$ during the forward pass leaves both the network’s overall output and the gradients of the loss with respect to the model parameters unchanged.

\paragraph{New Memory Mixing.} The sLSTM architecture, similar to the original LSTM, supports multiple memory cells (see Appendix~\ref{sec:appsLSTM}). The inclusion of multiple memory cells facilitates memory mixing through recurrent pathways, namely $\mathbf{R}_{\mathbf{z}}$, $\mathbf{R}_{\mathbf{i}}$, $\mathbf{R}_{\mathbf{f}}$, and $\mathbf{R}_{\mathbf{o}}$, which connect the hidden state vector $\mathbf{h}$ to the memory cell input $\mathbf{z}$ and the various gating mechanisms $\mathbf{i}$, $\mathbf{f}$, and $\mathbf{o}$. In this approach, exponential gating introduces a novel dimension to memory mixing. Furthermore, the redesigned sLSTM is capable of utilizing multiple heads, where memory mixing is permitted within individual heads but not between different heads. The combination of multi-head design in sLSTM and exponential gating paves the way for a distinct strategy of memory mixing.

\subsection{mLSTM}
\label{sec:mLSTM}

In order to improve the memory capabilities of LSTMs, we replace the conventional scalar memory cell $c \in \mathbb{R}$ with a matrix form $\mathbf{C} \in \mathbb{R}^{d \times d}$. This modification enables information retrieval through matrix multiplication. Specifically, at each time step $t$, we aim to encode a pair of vectors: the key $\mathbf{k}_t \in \mathbb{R}^d$ and the value $\mathbf{v}_t \in \mathbb{R}^d$ (adopting the terminology used in the Transformer architecture). Subsequently, at time $t + \tau$, a query vector $\mathbf{q}_{t+\tau} \in \mathbb{R}^d$ is used to access the stored value $\mathbf{v}_t$. This configuration aligns with the framework of Bidirectional Associative Memories (BAMs) \citep{Kohonen:72,Anderson:72,Nakano:72,Anderson:77}.

The covariance update rule~\citep{Sejnowski:77,Dayan:91} for storing a key-value pair is

\begin{align}
\mathbf{C}_t \ &= \ \mathbf{C}_{t-1} \ + \ \mathbf{v}_{t} \ \mathbf{k}_t^\top \ .
\end{align}

We posit that applying layer normalization prior to projecting inputs onto keys and values ensures that these components exhibit zero mean. The covariance update mechanism is known to be optimal~\citep{Dayan:91} for maximizing the distinction among retrieved binary vectors, which corresponds to achieving the highest possible signal-to-noise ratio. While greater separability can be obtained if retrieval is restricted to only pairwise relationships—albeit at the cost of incurring quadratic computational complexity, as seen in attention mechanisms~\citep{Krotov:16,Krotov:17,Ramsauer:21}—the covariance rule itself aligns with the formulation of Fast Weight Programmers \citep{Schmidhuber:92ncfastweights,Schlag:21}. Subsequent work introduced a constant decay factor applied to $\mathbf{C}_{t-1}$ and a fixed learning coefficient scaling 
$\mathbf{v}_{t} \mathbf{k}_t ^\top$~\citep{Ba:16}. Building on this perspective, we embed the covariance update approach within the LSTM model, mapping the forget gate to the decay parameter, the input gate to the learning coefficient, and using the output gate to modulate the vector retrieved from memory.

For this matrix memory, the normalizer state is the weighted sum of key vectors, where each key vector is weighted by the input gate and all future forget gates. Again, the normalizer state keeps record of the strength of the gates. Since the dot product between query and normalizer state can be close to zero, we use the absolute value of this dot product and lower bound it by a threshold (typically 1.0) as done previously~\citep{Sun:23arxiv}. The mLSTM forward pass is:

\begin{align}
\label{eq:mlstm_recurrent_begin}
\mathbf{C}_t \ &= \  f_t \ \mathbf{C}_{t-1} \ + \ 
  i_t \ \mathbf{v}_t \ \mathbf{k}_t^\top & &  &\text{cell state} \\
\mathbf{n}_t \ &= \  f_t \ \mathbf{n}_{t-1} \ + \ 
  i_t \ \mathbf{k}_t & &  &\text{normalizer state} \\
\mathbf{h}_t  \ &= \ \mathbf{o}_t \ \odot \ \tilde{\mathbf{h}}_t
\ , \qquad \qquad 
\tilde{\mathbf{h}}_t \ = \   \mathbf{C}_t \mathbf{q}_t \ / \ 
  \max \left\{ |\mathbf{n}_t^\top \mathbf{q}_t|, 1 \right\} 
   &&&\text{hidden state} \\
\mathbf{q}_t \ &= \ \mathbf{W}_q \ \mathbf{x}_t \ + \ \mathbf{b}_q  & &  &\text{query input} \\
\mathbf{k}_t \ &= \ \frac{1}{\sqrt{d}} \mathbf{W}_k \ \mathbf{x}_t \ + \ \mathbf{b}_k  & &  &\text{key input} \\
\mathbf{v}_t \ &= \ \mathbf{W}_v \ \mathbf{x}_t \ + \ \mathbf{b}_v  & &  &\text{value input} \\
i_t \ &= \ \!\exp \!  \! \left( \tilde{i}_t  \right) \ , 
 \qquad \qquad \ \ \ \ \,
  \tilde{i}_t \ = \ \mathbf{w}^\top_{i} \ \mathbf{x}_t \ + \  b_{i} \ \ \ 
  &&&\text{input gate} \\
f_t \ &= \ \sigma \!  \left(  \tilde{f}_t \right) \ \text{OR} \ \!\exp \!  \! \left(  \tilde{f}_t \right) , \ \ \: \!
\tilde{f}_t \ = \ \mathbf{w}^\top_{f} \ \mathbf{x}_t  \ + \
  b_{f}
  &&& \text{forget gate} \\
\label{eq:mlstm_recurrent_end}
\mathbf{o}_t \ &= \ \sigma \left( \tilde{\mathbf{o}}_t \right) \ , \qquad \qquad \quad \ \ \ \,
  \tilde{\mathbf{o}}_t  \ = \ \mathbf{W}_{\mathbf{o}} \ \mathbf{x}_t \ + \
  \mathbf{b}_{\mathbf{o}}
  &&&\text{output gate} 
\end{align}

mLSTM, akin to the standard LSTM, is capable of containing several memory cells. In the context of mLSTM, using multiple heads is functionally similar to employing multiple cells, as memory mixing is absent. To ensure the stability of the exponential gates in mLSTM, we apply identical stabilization methods to those utilized in sLSTM (refer to Equation~\ref{eq:slstmstabil}). Given the lack of memory mixing in mLSTM, its recurrence formulation allows for a parallelized version. Additional information can be found in Appendix~\ref{sec:appmLSTM}.

\subsection{xLSTM Architecture}
\label{sec:xLSTMarch}

\paragraph{xLSTM Blocks.}
\label{sec:xLSTMblock}
An xLSTM block is designed to encode historical information in a non-linear fashion within a high-dimensional space, thereby enhancing its capability to distinguish between diverse historical sequences or contextual backgrounds. Achieving such separation of histories is fundamental for accurately forecasting the subsequent element in a sequence, such as the next token. Our approach is motivated by Cover’s Theorem~\citep{Cover:65}, which posits that non-linear embedding of patterns into spaces of higher dimension increases the probability of achieving linear separability compared to the original, lower-dimensional space. We examine two types of residual block architectures: (i) The residual block employing post up-projection, as found in Transformer architectures, which first forms a non-linear summary of preceding information in the original dimension, then projects this to a higher-dimensional representation via a linear transformation, applies a non-linear activation, and finally maps it back to the original space; this design is illustrated in the left part of Figure~\ref{fig:Backbones} and the third column of Figure~\ref{fig:xlstm_sketch}. A more granular depiction can be found in Figure~\ref{fig:appsLSTM_detailed} within the appendix. (ii) The second, a residual block utilizing pre up-projection, similar in nature to those in State Space Models, initiates by mapping the inputs linearly into a higher-dimensional space,

The model captures historical information by applying a nonlinear transformation within a high-dimensional latent space, followed by a linear projection back into the original feature space. When an xLSTM block incorporates an sLSTM, we apply the up-projection block after the recurrent computation. Conversely, for an xLSTM block with an mLSTM, the up-projection is placed before the recurrent unit, capitalizing on the expanded memory capacity available in the higher-dimensional representation. Detailed schematic diagrams can be found in the left section of Figure~\ref{fig:Backbones}, the third column of Figure~\ref{fig:xlstm_sketch}, and Figure~\ref{fig:appsLSTM_detailed} within the appendix.

\paragraph{xLSTM Architecture.}
\label{sec:xLSTMarchitecure}
The xLSTM framework is assembled through the residual stacking of component blocks, following the methodology presented in~\citep{Srivastava:15,He:16}. This design employs the prevalent pre-LayerNorm~\citep{Ba:16b} residual structure, mirroring architectures found in state-of-the-art Large Language Models. Additional illustrations are available in the rightmost column of Figure~\ref{fig:xlstm_sketch}.

\subsection{Memory and Speed Considerations}

In contrast to Transformers, xLSTM architectures possess linear computational complexity and maintain fixed memory usage relative to sequence length. Owing to its compressive memory design, xLSTM is particularly advantageous for deployment in industrial contexts and for edge computing scenarios.

Although mLSTM does not use trainable parameters for its memory, it incurs high computational cost due to the $d \times d$ matrix memory and its corresponding updates of the same size. This approach balances an increased memory capacity with greater computational load. However, as these operations can be executed concurrently on GPUs, their overall impact on real-time execution remains minimal.

Unlike mLSTM—which, like FlashAttention~\citep{Dao:22,Dao:23} or GLA~\citep{Yang:23arxiv}, supports parallelization—sLSTM’s memory mixing (arising from hidden-to-hidden interactions) inhibits parallel computation. To address this, we have engineered an efficient CUDA implementation with optimized GPU memory usage down to the register level; as a result, sLSTM achieves execution speeds typically no more than twice as slow as mLSTM.

\section{Related Work}
\label{sec:related}

\paragraph{Linear Attention.} Numerous approaches have been proposed to address the issue of the Transformer's quadratic computational cost with respect to context length, aiming to achieve linear attention. The Synthesizer produces synthetic attention weights that are independent of pairwise token interactions~\citep{Tay:20arxiv}. Linformer approximates self-attention through a low-rank factorization of the attention matrix, yielding a linear operation~\citep{Wang:20arxiv}. The Linear Transformer reformulates the attention procedure to operate in linear time with context size~\citep{Katharopoulos:20}. Performer attains a linear approximation of the softmax attention via positive orthogonal random feature techniques~\citep{Choromanski:21}. Alternative frameworks such as the Structured Global Convolution (SGConv)~\citep{Li:22} and the Hyena Hierarchy~\citep{Poli:23} achieve comparable effects by substituting the attention layer with rapid, long-range convolutional operations.

\paragraph{State Space Models.} In recent years, State Space Models (SSMs) have gained significant traction due to their linear scalability with respect to sequence length and competitive performance relative to Transformers. The introduction of the Structured State Space sequence model (S4)~\citep{Gu:21} marked an early development in this area, soon followed by other architectures such as the Diagonal State Space (DSS) model~\citep{Gupta:22}, Gated State Space (GSS) models~\citep{Mehta:22}, the S5 model~\citep{Smith:22}, Bidirectional Gated SSM (BiGS)~\citep{Wang:22}, H3 model~\citep{Fu:23}, and Mamba~\citep{Gu:24arxiv}.

\paragraph{Recurrent Neural Networks.} Recurrent Neural Networks (RNNs) have been proposed as alternatives to Transformer architectures and attention mechanisms, primarily because their computational complexity scales linearly with sequence length. Models incorporating Deep Linear Recurrent Units (LRUs) have demonstrated strong performance in the domain of language modeling~\citep{Orvieto:23, De:24arxiv}, alongside architectures such as the Hierarchically Gated Linear RNN (HGRN)~\citep{Qin:23} and its successor HGRN2~\citep{Qin:24arxiv}. Among the most notable RNN-based strategies for large-scale language modeling is RWKV~\citep{Peng:23arxivshort, Peng:24arxivshort}, which delivers results on par with Transformer-based models.

\paragraph{Gating.}
Gating, central to the design of LSTM, has since been revisited and reimagined in numerous contemporary models. It features prominently in architectures such as HGRN~\citep{Qin:23}, HGRN2~\citep{Qin:24arxiv}, Gated Linear Attention (GLA)~\citep{Yang:23arxiv}, Gated State Space (GSS) models~\citep{Mehta:22}, Bidirectional Gated SSM (BiGS)~\citep{Wang:22}, Moving Average Equipped Gated Attention (MEGA)~\citep{Ma:22}, RWKV~\citep{Peng:23arxivshort}, and Mamba~\citep{Gu:24arxiv}.

\paragraph{Covariance Update Rule.}
To improve memory retention, we have integrated a matrix-based memory mechanism within the mLSTM cell that follows a covariance update rule. A similar philosophy underpins other architectures, including Fast Weight Programmers \citep{Schmidhuber:92ncfastweights,Schlag:21}, RWKV-5 and RWKV-6~\citep{Peng:24arxivshort}, Retention~\citep{Sun:23arxiv}, Linear Transformer~\citep{Katharopoulos:20}, and HGRN2~\citep{Qin:24arxiv}.

\paragraph{Most Related.} The models most comparable to xLSTM from a conceptual standpoint are Retention~\citep{Sun:23arxiv}, RWKV~\citep{Peng:23arxivshort,Peng:24arxivshort}, and HGRN2~\citep{Qin:24arxiv}. All these architectures make use of a matrix memory mechanism and/or utilize gating structures. Nevertheless, unlike the sLSTM proposed here, these models lack support for memory mixing. The ability to mix memory is crucial for effectively addressing state tracking challenges, rendering LSTMs inherently more expressive than both State Space Models (SSMs) and Transformers~\citep{Merrill:24,Deletang:23}. Such state tracking capabilities are indispensable for tasks like program evaluation or maintaining the continuity of entities throughout extended narrative sequences.

\paragraph{Residually Stacking Architectures.} As with the majority of modern large-scale neural networks, xLSTM models employ a design based on stacking modular components using residual connections~\citep{Srivastava:15,He:16}.

Such an approach facilitated the success of deep convolutional neural networks~\citep{He:16} and is foundational to the architecture of Transformers~\citep{Vaswani:17}. Transformers, in turn, underpin the development of Large Language Models (LLMs) including GPT-3~\citep{Brown:20short}, ChatGPT~\citep{Schulman:22short}, GPT-4~\citep{Achiam:23gpt4short}, Megatron-LM~\citep{Shoeybi:19}, Gopher~\citep{Rae:21short}, ERNIE 3.0 Titan~\citep{Wang:21arxivshort}, GLaM~\citep{Du:21glamshort}, the Chinese M6 model~\citep{Lin:21}, multilingual AlexaTM 20B~\citep{Soltan:22}, OPT~\citep{Zhang:22opt}, Chinchilla~\citep{Hoffmann:22short}, BLOOM~\citep{Scao:22arxivshort}, GLM-130B~\citep{Zeng:22arxivshort}, LaMDA~\citep{Thoppilan:22short}, PaLM~\citep{Chowdhery:22short}, Llama~\citep{Touvron:23arxiv}, and Gemini~\citep{GeminiTeam:23arxiv,Reid:24arxiv}.

\section{Experiments}
\label{sec:Exp}

This section presents a comprehensive empirical assessment of xLSTM, emphasizing its utility for language modeling tasks. Section~\ref{sec:ExpSyntethic} is devoted to analyzing xLSTM's properties on a series of synthetic benchmarks. Subsequently, in Section~\ref{sec:ExpComparison}, we report a comparative analysis of validation perplexities for leading language modeling approaches, each trained on a 15B-token sample from the SlimPajama dataset~\citep{Soboleva:23}. This analysis also includes ablation studies targeting components of xLSTM within the same training regime. In addition, we investigate how different models scale as a function of available data, following the methodologies established by \citet{Kaplan:20} and \citet{Brown:20short}.

Section~\ref{sec:ExpLanguage} expands the investigation by conducting extensive language modeling evaluations. Here, xLSTM is compared to top contenders identified in Section~\ref{sec:ExpComparison}, with all models trained on an expanded dataset of 300B tokens from SlimPajama~\citep{Soboleva:23}. The evaluation framework is multifaceted: (1) we analyze model performance on length generalization tasks; (2) validation perplexities are benchmarked alongside downstream task results as per \citet{Sutawika:24}; (3) evaluation is extended to 571 diverse text domains from the PALOMA language benchmark~\citep{Magnusson:23arxivshort}; (4) we reassess scaling dynamics among models, this time under conditions of significantly increased data volume—20-fold more than earlier experiments.

\label{sec:xlstm_blocks_notation}
Throughout all conducted experiments, we adopt the notation xLSTM[$a$:$b$] to denote a configuration in which the proportion of mLSTM-based blocks to sLSTM-based blocks is $a/b$. To illustrate, xLSTM[7:1] specifies an arrangement consisting of eight total blocks, with seven utilizing the mLSTM architecture and a single block employing the sLSTM architecture. In the case where the total number of blocks is fixed at 48, this corresponds to 42 mLSTM-based and 6 sLSTM-based blocks. Additionally, our experimental setup consistently utilizes up-projection blocks preceding mLSTM layers and following sLSTM layers.

\subsection{Synthetic Tasks and Long Range Arena}
\label{sec:ExpSyntethic}
We begin by evaluating the impact of xLSTM's novel exponential gating mechanism, combined with memory mixing, on tasks involving formal languages~\citep{Deletang:23}. Subsequently, we examine the performance of xLSTM's newly introduced matrix memory on the Multi-Query Associative Recall task~\citep{Arora:23arxiv}. Lastly, we investigate xLSTM's capabilities in handling extended sequences within the Long Range Arena benchmark~\citep{Tay:21}.

\paragraph{Evaluation of xLSTM's Exponential Gating with Memory Mixing.} We evaluate the newly introduced exponential gating mechanism with memory mixing in xLSTM, designed to address state tracking challenges~\citep{Merrill:24, Merrill:23}. Building on the formal language benchmarks of \citet{Deletang:23}, we modify and expand these tasks to support training on variable sequence lengths, allowing examination of the models' ability to extrapolate to longer sequences. A comprehensive outline of the tasks and additional experimental results are included in Appendix~\ref{sec:appExpSynthetic-formalLang}.  
Our assessment compares xLSTM with multiple alternatives such as Transformers, State Space Models, and classic Recurrent Neural Networks. Performance is measured as the proportion of relevant tokens classified correctly, normalized from 0 (chance-level) to 1 (fully accurate). We use two-block versions of each architecture for comparison: xLSTM[0:1] (sLSTM only), xLSTM[1:0] (mLSTM only), xLSTM[1:1], Llama, Mamba, RWKV, Retention, Hyena, traditional LSTM, and LSTM in Transformer blocks (LSTM (Block)). Figure~\ref{fig:formal-main} reports the empirical results.  
Notably, architectures lacking memory mixing and explicit state tracking, such as Transformers or State Space Models, are unable to handle certain regular grammar tasks, for instance the parity problem. This observation reinforces prior work demonstrating that both Transformers and State Space models are inherently less expressive for these tasks compared to RNN-based approaches~\citep{Merrill:24, Merrill:23, Deletang:23}.

\paragraph{Test of xLSTM's Memory Capacities on Associative Recall Tasks.} In this study, we examine the enhanced matrix memory mechanism of xLSTM within the context of the Multi-Query Associative Recall task~\citep{Arora:23arxiv}. Each sequence is comprised of randomly assigned key-value pairs drawn from an extensive vocabulary, requiring models to store these associations for subsequent retrieval. To rigorously assess memory capacity, we extend the original challenge by increasing both the number of key-value pairs—reaching up to 256—and the context window to a maximum of 2048 tokens. These more demanding conditions allow us to comprehensively evaluate memory performance across various models. Specifically, we benchmark 2-block variants of Llama, Mamba, RWKV-5, RWKV-6, xLSTM[1:1], and xLSTM[1:0], gauging their effectiveness through recall accuracy on the task. Transformers such as Llama, whose memory scaling is exponential with respect to the coding dimension~\citep{Ramsauer:21}, are used as the benchmark for optimal performance. Results can be found in Figure~\ref{fig:mqar-main}. Notably, xLSTM[1:1] demonstrates the highest memory performance among all non-Transformer baselines, including on models with fewer parameters. Remarkably, the inclusion of an sLSTM block appears not only to preserve but to strengthen memory capacity, an effect that is most pronounced on the hardest version of the task featuring 256 pairs. Further results, including extrapolation analyses suggesting that xLSTM's improved memory extends to sequence lengths beyond its training regime, are provided in Appendix~\ref{sec:appExpSynthetic-mqar}.

\paragraph{Test of xLSTM's Long Context Capabilities on Long Range Arena.} In order to evaluate how effectively xLSTM manages lengthy sequences and extensive contextual information, we benchmarked several approaches using the Long Range Arena~\citep{Tay:21}. Across all tasks, xLSTM achieved robust and consistently high results, indicating that the xLSTM architecture excels in efficiently addressing a variety of challenges associated with long context scenarios.

For more details, see Appendix~\ref{sec:appExpSynthetic-lra}.

\subsection{Method Comparison and Ablation Study}
\label{sec:ExpComparison}

In order to investigate the central inquiry of our study—specifically, the capabilities of our novel LSTM variants when applied at scale to language modeling tasks—we conduct experiments involving xLSTMs, Transformers, State Space Models, and additional approaches. All models are trained under identical auto-regressive conditions on a 15B token subset from SlimPajama. We subsequently assess these trained architectures on the validation dataset and carry out ablation analyses focused on the xLSTM models.

\paragraph{Comparative Evaluation of xLSTM and Baselines.} To facilitate comparison, we trained all models on a corpus of 15B tokens drawn from SlimPajama~\citep{Soboleva:23}. Performance assessment relies on perplexity measured on the validation split. The evaluated methods include our proposed xLSTM; two Transformer-based architectures, GPT-3~\citep{Brown:20short} and Llama~\citep{Touvron:23arxiv}; state space model architectures H3~\citep{Fu:23} and Mamba~\citep{Gu:24arxiv}; several recurrent neural network variants—RWKV-4~\citep{Peng:23arxivshort}, RWKV-5~\citep{Peng:24arxivshort}, RWKV-6~\citep{Peng:24arxivshort}, HGRN~\citep{Qin:23}, and HGRN2~\citep{Qin:24arxiv}—as well as GLA~\citep{Yang:23arxiv}, RetNet~\citep{Sun:23arxiv}, and Hyena~\citep{Poli:23}, which are all categorized as linear Transformers. Additionally, we include xLSTM[1:0] and xLSTM[7:1] variants (see Section~\ref{sec:xlstm_blocks_notation}). Mixed-precision training was utilized for all models, although for RWKV-5, RWKV-6, GLA, and HGRN2, automatic mixed precision via PyTorch was not used (additional information can be found in Appendix Section~\ref{sec:appGenTrainProc}). 

We organize the approaches under the following headings: (a) Transformers, (b) State Space Models (SSMs), and (c) Recurrent Neural Networks (RNNs), alongside linear Transformer architectures. Notably, linear Transformers implement a variant of the attention mechanism, substituting it with a linear operation. Each configuration was standardized to match the GPT-3 architecture comprising 350M parameters (specifically, embedding dimension of 1024 and 24 residual layers). GPT-3 is the sole approach to share weights between token and output embeddings, resulting in a parameter count advantage. As shown in Table~\ref{tab:spaj15b_model_results}, xLSTM achieves superior validation perplexity over all comparison models. Expanded results are reported in Appendix~\ref{sec:appExpComparison}. Figure~\ref{fig:temp_sclaw15B} illustrates the scaling trend observed, suggesting that xLSTM is likely to maintain its advantages as model size increases.

\paragraph{Ablation Studies.}
As shown in Table~\ref{tab:spaj15b_model_results} and Figure~\ref{fig:temp_sclaw15B}, xLSTM delivers outstanding language modeling performance after training on 15B tokens from SlimPajama. To systematically examine the contributions from LSTM to xLSTM, we incrementally convert a standard LSTM model into an xLSTM one. The process begins by embedding LSTM layers within pre-LayerNorm residual structures. Next, this configuration is expanded to include a post up-projection component. The final transformation introduces both exponential gating mechanisms and matrix memory capabilities.

The results are shown in Table~\ref{tab:ablstudies} (top). 

Ablation experiments credit the notable performance gains to the synergy between exponential gating and matrix memory. Given the crucial role of gating within both RNNs and State Space Models, we further investigate the effect of various gating strategies. As illustrated in Table~\ref{tab:ablstudies} (bottom), our findings demonstrate that learnable gates conditioned on the input systematically enhance model effectiveness. Further analysis concerning separate backbone elements is provided in Appendix~\ref{sec:appAblDetails}.

\subsection{xLSTM as Large Language Model}
\label{sec:ExpLanguage}

This work concludes with comprehensive experiments on large-scale language modeling, aimed at evaluating the viability of xLSTM as a large language model (LLM). For this purpose, we expand the training dataset to 300B tokens from SlimPajama—a scale consistent with prior work such as Mamba~\citep{Gu:24arxiv} and Griffin~\citep{De:24arxiv}.

xLSTM is benchmarked against RWKV-4, Llama, and Mamba, with each serving as a representative from the different methodological categories discussed in Section~\ref{sec:ExpComparison}. RWKV-4 is chosen to represent RNN-based methods, as suitable training precision configurations for RWKV-5, RWKV-6, and HGRN2 were only determined after initiating the 300B token experiments (see Appendix~\ref{sec:appGenTrainProc} for further details). 

Models of varying parameter counts (125M, 350M, 760M, 1.3B) are trained, and their capability for length extrapolation is tested alongside evaluation on a held-out validation set. We further analyze model performance on downstream tasks, assess language modeling effectiveness across the 571 text domains included in the PALOMA benchmark, and examine scaling law trends.

\paragraph{Sequence Length Extrapolation.}
\label{sec:spaj300B_ctx_extrapolation}
We begin by evaluating the capability of large-scale 1.3B parameter xLSTM, RWKV-4, Llama, and Mamba models to generalize to longer sequence lengths. Although training is performed with a fixed context length of 2048, we assess model performance for input lengths as large as 16384 tokens. The outcomes are summarized in Figure~\ref{fig:spaj300B_seq_extrapolation}. Notably, xLSTM consistently preserves low perplexity scores at extended context lengths, in contrast to the increased perplexity observed in other approaches.

\paragraph{Validation Perplexity and Downstream Tasks.}
\label{sec:spaj_downstream_eval}
Next, across all scales considered, we measure the results of xLSTM, RWKV-4, Llama, and Mamba on the SlimPajama validation split for the task of next token prediction as well as for a collection of downstream benchmarks focused on common sense reasoning capabilities.

The third column in Table~\ref{tab:lmeval} presents the validation set perplexities across the various approaches. For every model size, xLSTM[1:0] and xLSTM[7:1] consistently achieve the lowest validation set perplexity, establishing them as the leading models by this metric. The remaining columns in Table~\ref{tab:lmeval} display results for a range of downstream tasks. Across most tasks and model scales, xLSTM outperforms alternatives—Mamba only surpasses it in certain instances on the ARC task. Additional experimental results are provided in Appendix~\ref{sec:appExpLanguage}.

\paragraph{Performance on PALOMA Language Tasks.} Thirdly, we evaluate the next token prediction accuracy of the xLSTM, RWKV-4, Llama, and Mamba architectures across all model sizes, utilizing the PALOMA language tasks~\citep{Magnusson:23arxivshort}. Performance is quantified by measuring perplexity for next token prediction across 571 distinct textual domains, which encompass sources from nytimes.com to Reddit's r/depression. 

The results, summarized in Table~\ref{tab:ppleval}, report token prediction perplexities, categorizing them into broader language modeling tasks (first seven columns) and more specific domain benchmarks (final five columns). In this set of tasks, xLSTM[1:0] consistently yields lower perplexity scores compared to xLSTM[7:1].

Specifically, xLSTM[1:0] achieves a lower perplexity than Mamba on 568 of the 571 domains (99.5\%), outperforms Llama in 486 out of 571 cases (85.1\%), and surpasses RWKV-4 in 570 out of 571 domains (99.8\%). Additional details are provided in Appendix~\ref{sec:appExpLanguage}.

\paragraph{Scaling Laws.} Fourthly, we examine the power-law scaling characteristics, which enable predictions about performance as model size increases~\citep{Kaplan:20,Brown:20short}. Figure~\ref{fig:temp_sclaw300B} illustrates these scaling dynamics. Although all models exhibit comparable scaling patterns, there are distinct differences in their baselines. Among them, RWKV-4 demonstrates the least favorable performance, with Llama and Mamba following in ascending order. Notably, xLSTM outperforms Mamba, with the advantage xLSTM holds over Mamba roughly matching the margin between Mamba and Llama. The observed scaling trends suggest that as models grow, xLSTM is likely to maintain a relative performance edge over both Transformers and State-Space architectures.

\textbf{Generation Times and Maximal Throughput.} We present measurements of text generation latency (see Figure~\ref{fig:inference_time_generation}) and maximal achievable throughput (refer to Figure~\ref{fig:inference_time_decoding_throughput}, left) for our xLSTM variants at the 1.3B parameter scale. Our evaluations benchmark against comparably sized implementations of Mamba, Llama, and RWKV from HuggingFace, including a static key-value cache for Llama. During our assessment, both full cache compilation for the Transformer model and \texttt{torch.compile} compilation for Mamba were not supported and thus unavailable. All models were assessed under a batch size of 1 with a pre-fill value of 16 in the text generation scenarios, a setting chosen to optimally favor Transformer-based architectures.

Figure~\ref{fig:inference_time_generation} illustrates that both xLSTM and other recurrent baselines, such as Mamba and RWKV-4, demonstrate linear time complexity, in contrast to the quadratic scaling observed in Llama. For throughput analysis, we experimented with various batch sizes and prefill levels specifically for Llama.

Figure~\ref{fig:inference_time_decoding_throughput} (right) demonstrates that the xLSTM is able to support substantially larger batch sizes than Llama owing to its constant memory requirement, enabling it to attain superior throughput.

\section{Limitations}

(i) Unlike mLSTM, the memory interaction mechanism in sLSTM inherently prevents operations from being parallelized, which restricts the possibility of a highly parallel implementation. However, we have designed an efficient CUDA kernel for sLSTM, achieving a runtime that is under twice as slow as our parallel mLSTM implementation. (ii) At present, the CUDA kernels available for mLSTM are not tuned for speed; as a result, the current system operates at approximately one quarter of the speed of FlashAttention or the scan algorithm leveraged by Mamba. Performance closer to FlashAttention could be realized with more optimized CUDA kernels. (iii) The matrix-based memory component in mLSTM incurs significant computational demands, as processing involves $d \times d$ matrices. Nevertheless, because memory reading and updating do not require learned parameters and can rely on generic parallel matrix operations, the practical increase in computation time remains minimal. (iv) Attention must be given to appropriately initializing the forget gates. (v) As the matrix memory does not scale with sequence length, significantly extending context length may cause memory saturation for larger input sizes. However, empirical results indicate this does not pose a significant concern for sequence lengths up to 16k; further discussion can be found in Section~\ref{sec:spaj300B_ctx_extrapolation}. (vi) Owing to the high computational expense of conducting large-scale language experiments, neither architectural nor hyperparameter optimization—particularly for larger xLSTM models—was fully explored. Unlocking the complete capabilities of xLSTM will likely require substantial further tuning and optimization.

\section{Conclusion}
We have provided a partial response to the central inquiry: What are the outcomes of scaling LSTM architectures to billions of parameters in language modeling tasks? Our findings indicate that such scaling reaches at least parity with modern techniques including Transformers and State Space Models. By introducing exponential gating with memory mixing and a restructured memory design, we extended LSTM to form the xLSTM architecture. Comparative evaluations demonstrate that xLSTM achieves strong performance in language modeling, rivaling established methods such as Transformers and State Space Models. According to observed scaling laws, expanding xLSTM model sizes is likely to position them as strong contenders to existing Large Language Models based on Transformer architectures. Beyond language modeling, xLSTM exhibits promise for substantial contributions in areas like Reinforcement Learning, Time Series Prediction, and the simulation of physical phenomena.

\end{document}